\documentclass[11pt]{article}

%% ---------------------------------------------------------------
%% Supplementary Material
%% Sensitivity analysis: optimistic upper bound for the deep models
%% under test-fold epoch selection.
%% This file is a standalone supplement to the main manuscript
%% (paper_structure.tex) and is NOT included in the main document.
%% ---------------------------------------------------------------

\usepackage[margin=1in]{geometry}
\usepackage{graphicx}
\usepackage{booktabs}
\usepackage{multirow}
\usepackage{amsmath}
\usepackage[hidelinks]{hyperref}
\usepackage{parskip}
\usepackage{setspace}
\usepackage{caption}

\setstretch{1.15}
\setlength{\parindent}{0pt}
\setlength{\parskip}{6pt}
\captionsetup{font=small, labelfont=bf, width=0.95\linewidth}

%% Number supplementary tables/figures as S1, S2, ...
\renewcommand{\thetable}{S\arabic{table}}
\renewcommand{\thefigure}{S\arabic{figure}}
\renewcommand{\thesection}{S\arabic{section}}

\title{\textbf{Supplementary Material}\\[4pt]
Radiomics-Based Machine Learning and Deep Learning for Clinically Significant
Prostate Cancer Detection on Biparametric MRI}
\date{}

\begin{document}
\maketitle

%% ---------------------------------------------------------------
\section{Sensitivity analysis: optimistic upper bound for the deep models
under test-fold epoch selection}
\label{app:outerval}

A natural concern when classical and deep models are compared is whether the
deep architectures were disadvantaged by their training protocol. In the main
analysis, the number of training epochs for each deep model was selected on an
internal 80:20 split of the outer-training partition, after which the model was
refit on the complete outer-training partition; the held-out outer fold was
never used for any fitting or model-selection decision, exactly mirroring the
leakage-safe treatment of the classical ensembles. To bound how much the deep
models could gain from a more permissive---but optimistically biased---protocol,
we repeated the deep-model experiments while selecting the training epoch by
monitoring the held-out outer fold itself (up to 300 epochs). This procedure
deliberately leaks information from the evaluation fold into model selection and
therefore yields an \emph{upper bound} on deep-model performance rather than a
valid estimate of generalization; the classical models were left unchanged and
serve as the leakage-safe reference. All metrics are pooled out-of-fold values
with patient-level cluster-bootstrap 95\% confidence intervals, reported at the
fixed 0.5 probability threshold.

Even under this favorable protocol, the deep models did not overturn the main
conclusions (Table~\ref{tab:outerval}). In the radiomics-only setting---the
primary comparison---the leakage-safe Random Forest (AUROC 0.814) remained at
least as discriminative as every test-fold-selected deep model (best deep AUROC
0.794), so the classical ensemble retained its advantage despite the deep models
being granted a peek at the evaluation fold. Only in the combined conditions did
the optimistically tuned deep models edge ahead of the Random Forest, and then
only marginally and with broadly overlapping confidence intervals: by
$+0.009$ AUROC for the concatenation Transformer (0.837 vs.\ 0.828) and by
$+0.015$ AUROC for the dual-branch Transformer (0.843 vs.\ 0.828). These small,
leakage-inflated margins are consistent with the main finding that, on this
engineered radiomic feature space, deep tabular models do not confer a robust
discrimination advantage over a calibrated tree ensemble, and they reinforce
rather than weaken that conclusion: the deep models fail to clearly surpass the
classical baseline even when the evaluation protocol is tilted in their favor.

\begin{table}[htbp]
\centering
\caption{Optimistic upper bound for the deep models when the training epoch is
selected on the held-out outer fold (a deliberately leakage-prone protocol),
compared with the leakage-safe Random Forest reference carried over unchanged
from the main analysis. Threshold-free metrics on pooled out-of-fold predictions
with patient-level cluster-bootstrap 95\% confidence intervals; lower Brier is
better. Deep-model rows are optimistic upper bounds and are \emph{not} valid
estimates of generalization.}
\label{tab:outerval}
\resizebox{\textwidth}{!}{
\begin{tabular}{llccc}
\toprule
Condition & Model (protocol) & AUROC & AUPRC & Brier \\
\midrule
\multirow{4}{*}{Radiomics-only}
 & Random Forest (leakage-safe, reference) & 0.814 (0.790--0.837) & 0.605 (0.561--0.655) & 0.152 \\
 & Transformer (test-fold epoch) & 0.794 (0.766--0.819) & 0.585 (0.544--0.636) & 0.161 \\
 & CapsNet (test-fold epoch) & 0.761 (0.735--0.788) & 0.559 (0.519--0.608) & 0.202 \\
 & Transformer--CapsNet (test-fold epoch) & 0.793 (0.767--0.818) & 0.588 (0.545--0.634) & 0.177 \\
\midrule
\multirow{4}{*}{Radiomics$+$clinical (concat)}
 & Random Forest (leakage-safe, reference) & 0.828 (0.805--0.851) & 0.634 (0.591--0.683) & 0.147 \\
 & Transformer (test-fold epoch) & 0.837 (0.815--0.860) & 0.649 (0.606--0.697) & 0.154 \\
 & CapsNet (test-fold epoch) & 0.831 (0.809--0.853) & 0.646 (0.605--0.691) & 0.170 \\
 & Transformer--CapsNet (test-fold epoch) & 0.828 (0.806--0.852) & 0.652 (0.612--0.696) & 0.166 \\
\midrule
\multirow{4}{*}{Radiomics$+$clinical (dual)}
 & Random Forest (leakage-safe, reference) & 0.828 (0.805--0.851) & 0.634 (0.591--0.683) & 0.147 \\
 & Dual Transformer (test-fold epoch) & 0.843 (0.821--0.864) & 0.671 (0.628--0.716) & 0.154 \\
 & Dual CapsNet (test-fold epoch) & 0.840 (0.817--0.861) & 0.661 (0.620--0.706) & 0.151 \\
 & Dual Transformer--CapsNet (test-fold epoch) & 0.829 (0.806--0.851) & 0.621 (0.581--0.673) & 0.171 \\
\bottomrule
\end{tabular}
}
\end{table}

\end{document}
